%======================================================================
\section{모형}
\label{sec:model}

\begin{yoji}
같은 주소들 위에 허락 층과 송금 층을 놓고, 그 위에서 페이지랭크, C-PR, S-PR, 기준 점수를 정의해요. 명제 두 개가 섞는 비율 $\lambda$가 어디서 작용하는지, 그리고 깊이 1인 스펜더에서 허락 걷기가 무엇을 계산하는지 말해 줘요. C-PR은 두 층의 화살표를 모두 걸어서, 표~\ref{tab:bench}의 점수 가운데 푸는 시간이 가장 길었어요.
\end{yoji}

\subsection{한 주소 집합 위의 두 층}

$V=V_A\cup V_T$는 두 층 가운데 어느 하나에라도 있는 주소의 집합이에요(그림~\ref{fig:layers}). 영문 원고는 낱말을 이렇게 나눠 써요. 층과 그 화살표와 그 걷기에는 authorization(허락 층), 화살표를 만드는 사건에는 approval(허락하기), 허락한 금액에는 allowance(허락 금액)를 써요. 둘째 층도 층은 payment(송금 층), 그 사건은 transfer(송금)로 나눠요.

\paragraph{허락 화살표.}
기준 시점까지 주인 $u$가 스펜더 $v$에게 남긴 마지막 \texttt{Approval} 사건이 적어도 한 토큰에서 양수이면, 허락 화살표 $u\to v$가 있어요. 마지막 값이 0이면 그 쌍은 지워져요. 이것이 처리 과정이 실제로 재는 것이고, 지금 유효한 허락과는 세 가지가 달라요.
\begin{itemize}
\item 점수 기간이 시작되기 전에 정해지고 그 기간 안에 다시 나오지 않은 허락은 보이지 않아요.
\item \texttt{transferFrom}으로 일부를 꺼내 쓴 유한 허락도, 허락한 값을 그대로 가져요. 허락을 쓸 때 토큰이 \texttt{Approval} 사건을 꼭 남기지는 않기 때문이에요.
\item Uniswap의 Permit2 같은 허락 관리자에게 허락한 주인은, 그 관리자로 가는 화살표만 가져요.
\end{itemize}

\paragraph{송금 화살표.}
$u$가 $v$에게 양수 금액의 송금을 적어도 하나 보냈으면 송금 화살표 $u\to v$가 있어요. 이더(ETH)는 \texttt{Transfer} 로그 없이 움직이므로, 감싼 이더(WETH)로만 송금 층에 들어가요.

\paragraph{금액으로 단 무게.}
이 원고가 다루는 연산자는 두 층 모두에 금액으로 무게를 달아요.
\begin{equation}
\tilde w_A(u,v)=\sum_{\text{토큰}} g\bigl(a^{\star}(u,v,\text{토큰})\bigr),
\label{eq:layerA}
\end{equation}
\begin{equation}
\tilde w_T(u,v)=\sum_{e:\,u\to v} x_e\,\sigma(\Delta t_e).
\label{eq:layerT}
\end{equation}
\begin{itemize}
\item $a^{\star}$는 마지막 허락, $x_e$는 송금 $e$의 금액이에요.
\item $\sigma(\Delta t)=1/(1+\exp(k(\Delta t-t_0)))$는 나이 $\Delta t$(일)로 매기는 로지스틱 할인이에요. AWP가 쓰는 꼴이에요~\cite{do2023awp}. AWP 논문은 매개변수 값을 정하지 않아서, $k=0.01$과 $t_0=180$은 저자들이 고른 값이에요. 나이는 기준 시점에서 거꾸로 재요.
\item 허락 층은 할인하지 않아요.
\item 주 분석인 수정 분석(\ref{sec:revision}절)은 금액을 다섯 토큰의 달러 값으로 재고, 허락마다 상한을 둬요: $g(a)=\min(a,U)$, $U=19{,}000$달러. 계획은 원래 단위, 모든 ERC-20 토큰, $g(a)=a$로 정했어요.
\end{itemize}
$o_A(u)$와 $o_T(u)$는 각 층에서 나가는 무게의 합이에요. $P_A$와 $P_T$는 행마다 합이 1이 되게 나눈(행 정규화한) 행렬이고, 그 층에 나가는 화살표가 없는 점의 행은 0이에요.

\begin{figure}[!ht]
\centering
\begin{tikzpicture}[>=Stealth, font=\small,
  acc/.style={circle, draw=black!70, thick, minimum size=7mm, inner sep=0pt, fill=white},
  con/.style={rectangle, draw=black!80, thick, minimum size=6.8mm, inner sep=0pt, fill=black!12},
  auth/.style={->, thick, dashed, draw=authc},
  pay/.style={->, thick, draw=payc}]
  \node[acc] (u)  at (0,0)        {$u$};
  \node[acc] (w)  at (0,1.4)      {$w$};
  \node[acc] (x)  at (2.4,-1.4)   {$x$};
  \node[con] (r)  at (2.4,0.75)   {$R$};
  \node[con] (s)  at (4.8,-0.35)  {$S$};
  \node[acc, draw=black!40, text=black!60] (z) at (7.0,-0.35) {$z$};
  \draw[auth] (u) to[bend left=14] (r);
  \draw[pay]  (u) to[bend right=14] (r);
  \draw[auth] (u) to[bend right=10] (s);
  \draw[auth] (w) -- (r);
  \draw[auth] (x) -- (s);
  \draw[pay]  (r) -- (x);
  \draw[pay]  (s) -- (z);
  \draw[auth] (-0.4,2.75) -- (0.4,2.75) node[right, text=black] {허락(양수인 마지막 허락, 주인 $\to$ 스펜더)};
  \draw[pay]  (-0.4,2.3) -- (0.4,2.3) node[right, text=black] {송금(토큰 송금, 보낸 쪽 $\to$ 받는 쪽)};
  \node[con, minimum size=3.4mm] at (5.3,1.55) {};
  \node[right] at (5.5,1.55) {코호트 계약(점수를 받음)};
  \node[acc, minimum size=3.6mm] at (5.3,1.1) {};
  \node[right] at (5.5,1.1) {주인이나 상대방};
  \node[align=left, anchor=north west, inner sep=1pt] at (-0.45,-2.0)
    {$u\in B$: C-PR의 행 $=\lambda P_A(u,\cdot)+(1-\lambda)P_T(u,\cdot)$\\
     $w$, $x$: 나가는 화살표가 허락뿐이라 $\lambda>0$이면 $P_A$를 따라감\\
     $R$, $S$: 나가는 화살표가 송금뿐이라 $\lambda<1$이면 $P_T$를 따라감};
\end{tikzpicture}
\caption{같은 주소 집합 위의 두 층(개념도). 네모는 점수를 받는 코호트 계약이고, 동그라미는 사용자 계정 같은 주인과 상대방이에요. $z$는 코호트 계약이 송금했기 때문에만 나타나는 주소예요. 두 층 모두에 나가는 화살표가 있는 점은 $u$뿐이라, $0<\lambda<1$에서 $\lambda$에 따라 달라지는 행은 $u$의 행뿐이에요. $\lambda=1$이면 송금 화살표와 $z$가 빠지고, 계약 $R$과 $S$는 매달린 점이 돼요. $\lambda=0$이면 허락 화살표가 빠지고, $w$는 그래프를 떠나고 $x$는 매달린 점이 돼요. S-PR은 실선 화살표만 걷고, 뛸 때마다 허락 걷기의 점수에 비례해 다시 출발해요.}
\label{fig:layers}
\end{figure}

\subsection{다시 출발 분포가 있는 페이지랭크}

아래의 모든 점수는 다음 식의 고정점이에요.
\begin{equation}
\mathbf r^{(t+1)}=(1-d)\,\boldsymbol\pi+d\,P^{\top}\mathbf r^{(t)}+d\,\bigl(\mathbf 1_{\mathrm{dang}}^{\top}\mathbf r^{(t)}\bigr)\,\boldsymbol\pi.
\label{eq:pr}
\end{equation}
$P$는 이동 행렬, $\boldsymbol\pi$는 다시 출발 분포예요. $\mathbf 1_{\mathrm{dang}}$은 행이 모두 0인 점(매달린 점)을 가리키는 표시이고, 그 점의 몫은 $\boldsymbol\pi$를 따라 다시 나눠져요. $d=0.85$를 쓰고, 균등한 벡터에서 시작해서, $\ell_1$ 변화가 $10^{-8}$ 아래로 내려가거나 300번 반복하면 멈춰요.

\begin{puri}[식~\eqref{eq:pr}을 말로 읽으면]
걷는 사람은 한 걸음마다 확률 $d$로 지금 점의 화살표 하나를 따라가고(무게가 큰 화살표일수록 자주), 확률 $1-d$로 $\boldsymbol\pi$를 따라 아무 점으로나 뛰어요. 나가는 화살표가 없는 매달린 점에 도착하면 그 몫도 $\boldsymbol\pi$를 따라 다시 나눠져요. 이 걸음을 되풀이해서 더는 바뀌지 않게 된 값이 각 점의 점수예요.
\end{puri}

\subsection{결합 페이지랭크}

C-PR은 점마다 두 층의 행을 섞어요. 섞을 때는 그 점에 나가는 화살표가 있는 층만으로 다시 정규화해요.
\begin{equation}
P^{(\lambda)}_{u\cdot}=\alpha_A(u)\,P_{A,u\cdot}+\alpha_T(u)\,P_{T,u\cdot},
\label{eq:cpr}
\end{equation}
\begin{equation}
\alpha_A(u)=\frac{\lambda\,\mathbf 1[o_A(u)>0]}{D_\lambda(u)},\quad
\alpha_T(u)=\frac{(1-\lambda)\,\mathbf 1[o_T(u)>0]}{D_\lambda(u)},
\label{eq:alpha}
\end{equation}
여기서 $D_\lambda(u)=\lambda\,\mathbf 1[o_A(u)>0]+(1-\lambda)\,\mathbf 1[o_T(u)>0]$이에요.
\begin{itemize}
\item 한 층에만 나가는 화살표가 있는 점은 그 층을 무게 1로 따라가요.
\item 어느 층에도 나가는 화살표가 없는 점은 매달린 점이에요.
\item 다시 출발은 균등해요.
\item 주 설정 $\lambda=0.5$와 민감도 값 $0.25$, $0.75$는 계획에서 정했어요(\ref{sec:protocol}절).
\end{itemize}

$\lambda=1$이면 송금 층이 그 층에만 있는 점들과 함께 빠지고, 연산자는 허락 층만의 페이지랭크가 돼요. $\lambda=0$이면 송금 층만의 페이지랭크가 돼요. 이 원고는 이 둘을 허락 걷기와 송금 걷기라고 부르고, 두 규칙은 C-PR을 이 둘과 견줘요.

이 가족은 양 끝에서 끊겨 있어요. $\lambda$가 0과 1 사이에 있으면, 나가는 화살표가 모두 송금인 점(예를 들어 돈을 내보내기만 하고 아무에게도 허락하지 않는 계약)은 송금 층을 따라가요. 그런데 $\lambda=1$에서는 같은 점이 매달린 점이 되거나 그래프에서 사라져요. 거꾸로, 나가는 화살표가 모두 허락인 점은 안쪽의 모든 $\lambda$에서 허락 층을 따라가지만 $\lambda=0$에서는 매달리거나 사라져요. 작은 시험 자료로 돌리는 단위 시험이, 각 끝점이 한 층 걷기를 그대로 재현하는지 확인해요.

\begin{puri}[그림~\ref{fig:layers}로 보는 C-PR]
$\lambda=0.5$에서 $u$에 선 걷는 사람은 화살표를 따라갈 때 반은 허락 층의 행을, 반은 송금 층의 행을 써요. $w$와 $x$는 허락 화살표만 있어서 허락 층을, $R$과 $S$는 송금 화살표만 있어서 송금 층을 무게 1로 따라가요. 그래서 $\lambda$를 0.25에서 0.75로 바꿔도 행이 달라지는 점은 $u$뿐이에요. 반면 끝점으로 가면 그래프 자체가 바뀌어요. $\lambda=1$에서는 $R$과 $S$가 매달린 점이 되고, $\lambda=0$에서는 $w$가 사라지고 $x$가 매달린 점이 돼요.
\end{puri}

\subsection{씨앗을 쓰는 빼 보기 실험}

S-PR은 송금 걷기를 그대로 두고, 균등한 다시 출발 벡터만 허락 걷기의 점수 $\mathbf r_A$로 바꿔요. 이때 $\mathbf r_A$를 $V_T$로 줄이고 다시 정규화해요.
\begin{equation}
\pi_u=\frac{r_{A,u}}{\sum_{v\in V_T} r_{A,v}},\qquad u\in V_T,
\label{eq:spr}
\end{equation}
$u\notin V_A$이면 $r_{A,u}=0$이고, 분모가 0이 되면 균등 분포로 돌아가요. 그러면 허락은 걷는 사람이 뛰고 나서 어디에 내릴지를 정하고, 송금은 뛰기와 뛰기 사이에 어디로 갈지를 정해요. S-PR과 C-PR을 견주면 허락 화살표를 \emph{길}로 쓰는 효과와 \emph{씨앗}으로 쓰는 효과를 나눠 볼 수 있어요.

\subsection{기준 점수}

발표된 점수나 짝 원고의 점수 두 개와, 수 두 개를 기준으로 써요.
\begin{itemize}
\item \textbf{AWP}~\cite{do2023awp}는 송금 층을 걸어요. 무게는 송금마다 $\sigma(\Delta t_e)V(x_e)$를 더한 것이고, $V(z)=2/(1+e^{-bz})-1$은 상한이 있는 값 변환이에요. 각 주소에서 최근 송금 활동에 비례해 다시 출발해요.
\item \textbf{엔도스랭크}~\cite{kwon2026companion}는 허락 층을 같은 상한 있는 무게로 걸어요. 무게는 마지막 허락과 그 나이에 매기고, 다시 출발은 균등해요.
\item \textbf{받은 허락 수}: 그 계약에 양수의 마지막 허락을 가진 서로 다른 주인의 수예요.
\item \textbf{송금 입차수}: 자기 말고 그 계약에 송금을 보낸 서로 다른 주소의 수예요.
\end{itemize}
두 수는 모두 동결 시점에 세요.

수정 분석에서는 $b=\ln 3/m$이에요. $m$은 2026년 3월 31일까지 12개월의 중앙 금액을 유효숫자 두 자리로 반올림한 값이에요. $b=\ln 3/m$이므로 금액 $m$의 무게가 2분의 1이 돼요. 송금은 $m_T=170$달러(169.3에서), 유한 허락은 $m_A=100$달러(99.97에서)예요. 계획에서는 원래 단위 위에서 $b=1$이었고, 그러면 원래 단위로 10개 이상인 금액은 모두 $V(z)$와 1의 차이가 $10^{-4}$보다 작아요. 두 기준 걷기는 사건마다 무게에 상한을 두지만, C-PR은 상한까지는 금액이 제한 없이 무게에 들어가요.

\begin{puri}[상한이 있는 무게 $V$는 어떻게 움직이나요?]
$V$는 금액이 0이면 0이고, 금액이 커질수록 1에 다가가지만 1을 넘지 않아요. $b=\ln 3/m$으로 두면 중앙 금액 $m$의 무게가 정확히 2분의 1이고, 그 두 배 금액은 \ill{0.8}, 절반 금액은 \ill{0.27}쯤이에요. 그래서 몇백 달러짜리 허락도 무제한 허락과 거의 같은 무게를 가져요. 반면 C-PR의 허락 층은 상한 $U$까지 금액에 비례해 무게를 주고, 송금 층은 상한 없이 금액에 시간 무게를 곱해요.
\end{puri}

\subsection{무게가 작용하는 곳}

$B=\{u: o_A(u)>0 \text{ 이고 } o_T(u)>0\}$를 두 층 모두에 나가는 화살표가 있는 점의 집합이라고 해요.

\begin{proposition}
\label{prop:localize}
$\lambda\in(0,1)$이고 $\boldsymbol\pi$가 균등하면, $B$ 밖의 $P^{(\lambda)}$ 행은 $\lambda$에 따라 달라지지 않아요. 그리고 식~\eqref{eq:pr}의 정상 벡터를 $\lambda$로 미분한 $\mathbf r'(\lambda)$는 다음을 만족해요.
\begin{equation}
\mathbf r'(\lambda)=d\,\bigl(I-d\,G_\lambda^{\top}\bigr)^{-1}\sum_{u\in B} r_u\bigl(\mathbf p_{A,u}-\mathbf p_{T,u}\bigr).
\label{eq:sens}
\end{equation}
여기서 $\mathbf p_{A,u}$와 $\mathbf p_{T,u}$는 $P_A$와 $P_T$의 행을 세운(전치한) 벡터이고, $G_\lambda$는 $P^{(\lambda)}$에서 매달린 점의 행을 $\boldsymbol\pi^{\top}$로 바꾼 행렬이에요. 따라서 $r(B)=\sum_{u\in B}r_u$로 두면 $\|\mathbf r'(\lambda)\|_1\le 2d\,r(B)/(1-d)$예요.
\end{proposition}

\begin{proofko}
$B$ 밖에서는 식~\eqref{eq:alpha}의 표시 함수 가운데 하나가 0이에요. 그래서 그 행은 $\lambda$가 얼마든 한 층의 행이거나 0이고, $(0,1)$에서 매달린 점의 집합도 바뀌지 않아요. $u\in B$이면 $D_\lambda(u)=1$이고 행은 $\lambda P_{A,u\cdot}+(1-\lambda)P_{T,u\cdot}$이므로, 그 미분은 $P_{A,u\cdot}-P_{T,u\cdot}$예요. $\mathbf r$의 합이 1이므로 식~\eqref{eq:pr}은 $\mathbf r=(1-d)\boldsymbol\pi+d\,G_\lambda^{\top}\mathbf r$로 쓸 수 있고, 이것을 미분하면 식~\eqref{eq:sens}이 나와요. $G_\lambda^{\top}$는 열의 합이 1인 행렬이라 $\|(I-dG_\lambda^{\top})^{-1}\|_1\le 1/(1-d)$이고, 확률 행 두 개의 차이는 $\ell_1$ 노름이 많아야 2예요.
\end{proofko}

이 명제는 무게가 \emph{어디서} 작용할 수 있는지를 말할 뿐, 그 작용이 작다고 말하지는 않아요. $r(B)$가 크면 이 한계는 두 분포 사이 $\ell_1$ 거리의 뻔한 한계인 2를 넘어 버려요. 이 원고의 그래프에서는 허락 층의 거의 모든 점이 나가는 화살표를 가진 주인이에요. 그래서 추출한 송금의 보낸 쪽으로도 나오는 주인은 모두 $B$에 들어요. 저자들은 $r(B)$를 계산하지 않았어요. 부록~\ref{app:explore}가 안쪽 $\lambda$ 값들 사이에서 순위가 얼마나 움직이는지를 보고해요.

\begin{puri}[명제~\ref{prop:localize}을 쉽게 말하면]
$\lambda$가 0과 1 사이에 있는 동안에는, 나가는 화살표가 한 층에만 있는 점은 $\lambda$를 바꿔도 따라가는 행이 그대로예요. $\lambda$가 바꾸는 것은 두 층 모두에 화살표가 있는 점들($B$)의 행뿐이에요. 그래서 점수가 $\lambda$에 따라 움직이는 크기는 걷는 사람이 $B$에 머무는 몫 $r(B)$에 묶여 있어요. 다만 이 그래프에서 그 몫이 작은지는 재지 않았어요.
\end{puri}

\subsection{깊이 1인 스펜더에서의 허락 걷기}
\label{sec:closed}

이 원고의 그래프는 한쪽 끝에 코호트 계약이 있는 사건만 담아요. 그래서 허락 층의 점은 대부분 허락을 주기만 하거나(순수한 주인) 받기만 해요(순수한 스펜더). 짝 연구가 이 구조에서의 닫힌 식(반복 계산 없이 바로 쓰는 식)을 끌어냈고~\cite{kwon2026companion}, 이 원고는 그 식을 증명과 함께 다시 적어요.

\begin{proposition}
\label{prop:closed}
식~\eqref{eq:pr}에서 $\boldsymbol\pi$가 균등하면, 들어오는 화살표가 없는 점은 모두 같은 점수 $c$를 가져요. 그리고 들어오는 이웃이 모두 들어오는 화살표가 없는 점 $v$의 점수는 $r_v=c\,(1+d\,\delta(v))$예요. 여기서 $\delta(v)=\sum_u w(u,v)/o(u)$는 $v$가 받은 무게를 주인 쪽에서 정규화한 합이에요.
\end{proposition}

\begin{proofko}
다시 출발이 균등하고 매달린 점의 몫이 $r_D$이면, 식~\eqref{eq:pr}은 $r_v=(1-d+d\,r_D)/|V|+d\sum_u r_u\,w(u,v)/o(u)$로 읽혀요. 들어오는 화살표가 없는 점은 첫째 항만 받고, 이것을 $c$라고 불러요. $v$의 들어오는 이웃 $u$가 모두 들어오는 화살표가 없으면 $r_u=c$이고, 이것을 넣으면 $r_v=c+d\,c\,\delta(v)$가 나와요.
\end{proofko}

이 명제는 어떤 화살표 무게에서도 성립해요. 그래서 엔도스랭크와 $\lambda=1$의 C-PR 모두에, 각자의 $\delta$로 들어맞아요.
\begin{itemize}
\item 수정 분석: 2026년 6월 30일 동결의 엔도스랭크 그래프에서, 코호트 스펜더 6,224개 가운데 5,232개가 깊이 1이에요. 그 모두에서 계산한 점수가 닫힌 식과 1\% 안으로 맞아요. 모든 스펜더로 넓히면 87.5\%가, 더 깊은 스펜더만 보면 21.9\%가 1\% 안이에요~\cite{kwon2026companion}.
\item 등록 분석: 26,512개 가운데 20,942개가 깊이 1이고 그 모두가 1\% 안이며, 모든 스펜더로는 86.4\%, 더 깊은 스펜더로는 35.1\%예요.
\end{itemize}
따라서 대부분의 스펜더에서 허락 걷기는 허락한 주인 수에 무게를 다시 단 것이에요. C-PR도 허락 층을 따라가는 곳에서는 이 한계를 그대로 물려받아요.

\subsection{비용}
\label{sec:cost}

식~\eqref{eq:pr}의 반복 한 번에는 $O(|E|)$가 들어요~\cite{leskovec2020mining}. C-PR은 두 층 화살표의 합집합을 걸어요. 표~\ref{tab:bench}은 2026년 6월 30일 동결 시점의 전체 그래프에서 각 점수의 풀이 시간을, 각 풀이 그래프의 $|V|$, $|E|$와 함께 보여 줘요.
\begin{itemize}
\item 수정 분석: C-PR은 점 4,111,490개와 화살표 10,562,390개에서 11.56초가 걸려요. AWP의 1.25배인 화살표에서 AWP 시간의 1.36배예요.
\item 등록 분석: 점 25,304,106개와 화살표 71,375,442개에서 105.51초가 걸려요. AWP의 1.19배인 화살표에서 AWP 시간의 1.51배예요.
\end{itemize}
엔도스랭크는 그래프가 가장 듬성듬성해서 가장 빨라요. 허락 그래프의 화살표는 송금 그래프의 45.2\%(수정 분석), 37.4\%(등록 분석)예요. AWP와 엔도스랭크의 시간 비 3.21과 2.36은, 이 화살표 비(2.21과 2.68)와 표~\ref{tab:bench}의 반복 수가 합쳐진 결과예요. AWP는 수정 분석에서 94번, 엔도스랭크는 57번 반복하고, 등록 분석에서는 93번과 84번 반복해요. 그러니 속도 차이는 화살표 수와 그 그래프에 필요한 반복 수를 반영할 뿐, 알고리즘의 차이가 아니에요. 송금 층을 쓰는 걷기는 무엇이든 그 화살표의 값을 치러요. 시간은 한 기계(Intel Core Ultra 9 185H, 메모리 64\,GB, Windows 11, Python 3.11.9, NumPy 2.3.2, SciPy 1.16.1)에서, 시간을 재지 않는 준비 실행 한 번 뒤에 잰 다섯 번의 평균이에요.

\begin{table}[!ht]
\centering
\caption{2026년 6월 30일 동결 시점의 전체 그래프에서 잰 풀이 시간(준비 실행 한 번 뒤 시간을 잰 다섯 번의 평균과 표준편차)}
\label{tab:bench}
\small
\begin{tabular}{@{}lrrrr@{}}
\toprule
점수 & $|V|$ & $|E|$ & 반복 수 & 시간(초) \\
\midrule
\multicolumn{5}{@{}l}{\emph{수정 분석}} \\
EndorseRank & 2{,}513{,}982 & 3{,}808{,}949 & 57 & 2.66 (0.11) \\
AWP & 3{,}601{,}727 & 8{,}426{,}124 & 94 & 8.53 (0.62) \\
C-PR ($\lambda=0.5$) & 4{,}111{,}490 & 10{,}562{,}390 & 82 & 11.56 (0.21) \\
S-PR & 3{,}603{,}300 & 8{,}432{,}635 & 151 & 10.88 (0.33) \\
\midrule
\multicolumn{5}{@{}l}{\emph{등록 분석}} \\
EndorseRank & 13{,}055{,}150 & 22{,}321{,}470 & 84 & 29.60 (0.61) \\
AWP & 23{,}029{,}298 & 59{,}737{,}617 & 93 & 69.82 (2.30) \\
C-PR ($\lambda=0.5$) & 25{,}304{,}106 & 71{,}375{,}442 & 92 & 105.51 (2.48) \\
S-PR & 23{,}029{,}298 & 59{,}737{,}617 & 185 & 100.86 (4.82) \\
\bottomrule
\end{tabular}

\vspace{4pt}
\parbox{0.92\linewidth}{\footnotesize 그래프와 설정마다 준비 실행 한 번과 시간을 잰 실행 다섯 번을 한 차례 했고, 본문이 인용하는 측정이 이것이에요. 시간은 미리 만든 화살표 목록으로 행렬을 짜고 거듭제곱 반복을 하는 데 걸린 시간이에요($d=0.85$, 허용 오차 $10^{-8}$, 반복은 많아야 300번). S-PR의 시간에는 다시 출발 벡터를 주는 허락 층의 풀이도 들어 있어요. $|V|$와 $|E|$는 각 풀이 그래프의 값이에요. 풀이 그래프는 무게가 부동소수점에서 0이 되는 쌍을 빼므로, 수정 분석에서는 엔도스랭크와 AWP의 그래프에서 먼지 금액 쌍이 빠져요(\ref{sec:graphs}절). 두 분석 모두 같은 기계에서 쟀어요(\ref{sec:cost}절).}
\end{table}
